\section{The collision experiment and the results}
\label{sec:two-field-introduction}
The unknown launch position creates a joint inverse problem: one must
recover both the obstacle geometry and the law that blurs each nominal
command. We prove that two fixed opposite forward commands suffice.
Each is observed as a spatial field, and both use the same unknown
stationary law. Their difference recovers occupation by a finite
prefix formula. Their geometric supports then separate the obstacle
from the footprint by an exact cancellation. This gives a complete
inverse, including the probability law and the full translation
period group. The finite construction uses the same two commands.

\subsection{Observations and the exact inverse}
Let $\mathcal O$ be the union of closed obstacles. For a displacement
$b$, the bit $B_b(y)$ is one when $y\notin\mathcal O$ and the
unreflected segment $[y,y+b]$ meets $\mathcal O$. It is zero for a
solid start and for a free miss. A boundary solid start is zero; a
boundary contact reached from a free start counts as a collision.
For nominal center $x$ the hidden start is $x+Z$, where $Z$ has one
unknown law $\mu$. The observed mean is
\[
                 F_b(x)=\int B_b(x+z)\,\dd\mu(z).
\]
All attempted preparations enter the denominator. No free point,
actual start, obstacle label, contact position, angle or time is
observed. Fresh preparations are independent, conditional on the
chosen commands and prior observations. The finite theorems use exact
nominal controls and count their required numerical precision.

Fix $a=te$, with $|e|=1$ and $t>0$. The exact datum is the ordered
pair of whole-plane spatial fields $(F_a,F_{-a})$, supplied almost
everywhere. It consists of two command settings, not two scalar
measurements. Theorem~\ref{thm:two-field-rigidity} applies to nonempty
locally finite configurations of strictly convex bodies with bounded
diameters and a positive separation $d$. Their boundaries need not
be smooth. The law is any compactly supported probability whose
support $A$ is convex, including a point or segment, and
$t+\diam A<d$. The theorem reconstructs
\[
       \mathcal O-s(A),\qquad A-s(A),\qquad
                   (z\mapsto z-s(A))_\#\mu,
\]
Here $n_\phi=(\cos\phi,\sin\phi)$ and the Steiner point is
$s(K)=\pi^{-1}\int_0^{2\pi}h_K(n_\phi)n_\phi\dd\phi$. This common translation is the
entire ambiguity. Corollary~\ref{cor:two-field-fiber} identifies the
full period group and completes every finite-length forward response.
Neither exact conclusion assumes periodicity.

The central identity uses a matched occupation component $P=C-A$
and the two corresponding collision supports $K_+,K_-$. It is
\[
  2h_P(u)-h_{K_+}(u)-h_{K_-}(u)+t|u\cdot e|=h_C(u)-h_{L_C}(u),
\]
where $L_C$ joins the two transverse contact points of $C$. The
unknown support $h_{-A}$ cancels. The angular curvature measure on
the right consists of a nonatomic positive measure and the two
negative atoms of a segment. Its Jordan decomposition determines
$C$ up to translation. This uses the classical planar support-measure
calculus \cite{Schneider,MartinezMaure}; the observable cancellation
and its use for the joint collision inverse are the present argument.
The subsequent division of a compact occupation transform identifies
the law exactly. The finite law theorem instead uses moment estimates
with explicit conditioning.

These two conclusions have distinct quantitative content. The exact
theorem identifies a characteristic function by continuity from a
dense set of frequencies; it requires no lower bound for the obstacle
transform near its zeros. Finite acquisition first estimates geometric
supports, then fits a positive probability with controlled finitely
many moments. Its transportation error and its exponential sufficient
observation cost are proved directly. Exact response completion and
finite prediction are stated separately in their respective data
models and norms.

\subsection{Quantitative classes and losses}
For finite geometric reconstruction put $s=6+\beta$, where
$0<\beta\le1$. A bounded periodic presentation is
\begin{equation}\label{eq:two-field-presentation}
 \mathcal O=\bigcup_{i=1}^r(C_i+L\Z^2),\qquad 1\le r\le r_0.
\end{equation}
The presentation has exactly $r$ component orbits and may be
nonprimitive. Fix known constants with
\begin{equation}\label{eq:two-field-geometric-priors}
 \begin{gathered}
 \|L\|,\|L^{-1}\|\le L_0,\quad |\det L|\ge V_0>0,\quad
 \diam C_i\le D_0,\quad \dist(C,D)\ge d_0>0,\\
 \kappa_-\le\kappa_C\le\kappa_+,\qquad
                 \|p_C^\circ\|_{C^{6,\beta}}\le M_0.
 \end{gathered}
\end{equation}
Here distinct $C,D$ are physical components, $h_C$ is the support
function, $z_C=s(C)$, and $p_C^\circ=h_C-z_C\cdot n_\phi$.
The positive curvature bounds are uniform. The footprint and its
density have the separate priors
\eqref{eq:two-field-finite-priors}: a smooth support with upper and
lower curvature radii and a lower boundary-mass bound. These
quantitative hypotheses are used only in the finite assertions.

We state the finite period margin explicitly. Put $B=1+L_0$,
$Q_R=[-R,R]^2$, and $\Pi=\operatorname{Per}(\mathcal O)$. Define
\begin{align}
 d_*(C,D)&=\|z_C-z_D\|_\infty+
                         \|p_C^\circ-p_D^\circ\|_\infty,\nonumber\\
 \mathcal D(w)&=\max_{C:z_C\in Q_B}\ \max_{\sigma\in\{-1,1\}}
                            \inf_D d_*(C+\sigma w,D).
                   \label{eq:two-field-patch-defect}
\end{align}
The class $\mathcal B_\eta$ consists of these presentations with a
known $\eta>0$ such that
\begin{equation}\label{eq:two-field-patch-margin}
 \begin{gathered}
 w=z_D-z_C,\quad z_C\in Q_{3B},\ z_D\in Q_{5B},\quad w\notin\Pi
                 \quad\Longrightarrow\quad\mathcal D(w)\ge\eta.
 \end{gathered}
\end{equation}
This is a defect of whole-patch translations. It is invariant under
the common centering, as proved in the companion,
Lemma~\ref{H-lem:registration-period-margin}. The geometric loss is
the maximum $C^2$ error of matched support functions in a protected
patch, together with the errors of a matched primitive period basis
and free area. A basis is matched to a true basis; no continuous
choice of a reduced basis is imposed. Primitive orbit counts and
exact rational relations of the recovered period-generating list are
separate discrete outputs. For a finite configuration the loss is
the maximum $C^2$ error of all matched components in its known
complete protected aperture, and no period margin is assumed.

Theorem~\ref{thm:two-field-finite-geometry} recovers geometry to
$O(\nu)$ from finitely many bits with the fixed commands $\pm te_1$,
where $2t+\Delta<d_0$. Its sufficient power is
\[
            Q_{\rm pair}=\frac{(\gamma+9/2)s}{s-2}.
\]
Theorem~\ref{thm:two-field-finite-law} adds recovery of the canonical
probability in $W_1$ at sufficient cost
\[
 \exp(C\varepsilon^{-1}\log(C/\varepsilon))\log^2(C/\delta).
\]
On a protected complete component, Theorem~\ref{thm:linear-moment-sampling}
acquires all moments through degree $n$ with tolerance $\tau$ using
$C\tau^{-2}\log(C(n+1)/\delta)$ attempted bits. The observations are
shared across the moments, and no density regularity enters this
conditional acquisition bound. Corollary~\ref{cor:linear-joint-budget}
incorporates it into the joint inverse, keeping the geometric cost
and moment conditioning explicit.
Corollary~\ref{cor:two-field-finite-prediction} gives the corresponding
local spatial $L^1$ prediction of every bounded-length command.
Under a known total-variation bound on the zero-extended density,
Corollary~\ref{cor:two-field-finite-density} gives density $L^1$ error
and uniform raw-mean error on bounded windows. Its sufficient cost
has $\varepsilon^{-2}$ in place of $\varepsilon^{-1}$. The proofs
give these norms and costs directly; exact Fourier identification is
not used as a finite stability estimate.

\subsection{Resources and relation to the retained results}
We distinguish the number $S$ of distinct nominal sites, the number
$J$ of site-and-sign command occurrences, and the number $N$ of
attempted bits, including repetitions, solid starts and free misses.
Thus $S\le J\le N$. The two displacement vectors are fixed; a
required nominal grid is part of the theorem. Controller description
length, deterministic numerical arithmetic, and physical motion or
positioning are separate resources. A fixed nominal aperture becomes
an absolute physical aperture after a coarse bound on the unknown
footprint origin is included. This bound does not register the origin.

The technical supplement retains the directional-germ inverse, its
one-resolved-obstacle extension, all earlier preparation experiments,
and their complete proofs. The new fixed-pair inverse removes both
the angular continuum of commanded directions and the short-length
limit. Its geometric hypotheses do not compare the footprint diameter
with an obstacle width. For the quantitative smooth class, the
sufficient exponent $Q_{\rm pair}$ is one third of the retained
single-law exponent $(3\gamma+27/2)s/(s-2)$. At $\gamma=0$, $s=7$,
these powers are $6.3$ and $18.9$, respectively. These are sufficient
upper bounds for their stated experiments. The sharp power $2.3$
at $s=7$ in the companion concerns the separate fixed, known
uniform-disk law; it supplies no matched lower bound for the present
unknown-law problem.

The geometric hull step uses known nominal centers with positive
bits and a cap-mass argument of the kind underlying random support
estimation \cite{BrunelCaps}. It does not observe a random point of
the unknown obstacle. The law step uses classical moment comparison
and transportation duality \cite{RigolletWeed}, with a recovered and
therefore perturbed convolution factor. The prediction step uses
bounded variation \cite{AFP,Galerne}. Each of these reductions and
its uniform estimates is proved below. The supplement supplies the
radial interpolation and period-locking proofs through the explicit
interface in Section~\ref{sec:finite-interface}. The three main
sections proceed from exact identification, through finite geometry,
to finite law recovery and prediction. Section~\ref{sec:linear-moment-acquisition}
then sharpens the separated acquisition budget by a linear prefix rule. Appendix~\ref{sec:two-field-geometric-details}
gives the complete collar, grazing-cap, grid and assignment estimates;
Appendix~\ref{sec:moment-factor-details} gives the positive-factor
and finite probability reconstruction details.
